\chapter{Introducción}
\comentario{Este capítulo presenta la relevancia clínica de \textit{A. baumannii}, fundamenta la necesidad de integrar evidencia para priorizar dianas terapéuticas y establece los objetivos de la tesis.}

La resistencia a los antimicrobianos compromete la capacidad de tratar infecciones bacterianas y limita las opciones terapéuticas disponibles. En este contexto, la Organización Mundial de la Salud (OMS) incluyó a \textit{Acinetobacter baumannii} resistente a carbapenémicos entre los patógenos bacterianos de prioridad crítica, una categoría que busca orientar la investigación y el desarrollo de nuevos antibacterianos \cite{WHO2024BPPL}. La persistencia de este problema, junto con una cartera limitada de antibióticos que actúen sobre mecanismos novedosos, subraya la necesidad de identificar dianas terapéuticas que permitan ampliar las estrategias de tratamiento.

\textit{A. baumannii} es un bacilo gramnegativo, aerobio y no fermentador que se ha consolidado como un patógeno oportunista asociado principalmente a infecciones adquiridas en el ámbito hospitalario. Puede causar neumonía asociada a ventilación mecánica, bacteriemia, infecciones urinarias, de heridas y del sistema nervioso central, con especial impacto en personas críticamente enfermas o inmunocomprometidas \cite{Peleg2008}. Su relevancia clínica se vincula con la capacidad de adquirir y combinar determinantes de resistencia, modificar la permeabilidad de la envoltura celular, aumentar el eflujo de antimicrobianos, inactivar fármacos y alterar sus dianas \cite{Zhu2024}. Estas propiedades dificultan tanto el tratamiento de las infecciones como la selección de nuevos antibacterianos con actividad sostenida.

\section{Identificación de dianas terapéuticas en \textit{A. baumannii}}

El descubrimiento de una diana antibacteriana comienza por reconocer procesos cuya perturbación reduzca la viabilidad, el crecimiento o la capacidad de causar enfermedad del microorganismo. Sin embargo, la relevancia biológica de un gen no garantiza que su producto sea una diana terapéutica útil. Una candidata debe, además, estar suficientemente conservada entre los aislados que se busca tratar, diferenciarse de las proteínas del hospedador para reducir posibles efectos sobre proteínas no objetivo y presentar propiedades compatibles con la unión selectiva de una molécula inhibidora. Por ello, la selección de dianas requiere integrar evidencia funcional, evolutiva, anotacional y estructural, en lugar de depender de un único criterio \cite{Bosch2024}.

La diversidad genómica de \textit{A. baumannii} hace especialmente necesaria esta perspectiva integradora. La especie posee un repertorio génico variable, moldeado por transferencia horizontal, elementos móviles y recombinación, procesos que favorecen su adaptación y la circulación de determinantes de resistencia \cite{Touchon2014}. En consecuencia, una diana identificada en una cepa de referencia puede estar ausente, ser divergente o no desempeñar el mismo papel en otros aislados clínicos. El análisis de un pangenoma permite distinguir genes ampliamente conservados de aquellos restringidos a linajes particulares y, por lo tanto, priorizar candidatos con mayor potencial de cobertura dentro de la diversidad clínica de la especie.

\section{Esencialidad genética como evidencia funcional}

La esencialidad genética aporta una fuente de evidencia central para la priorización de dianas antibacterianas. En términos operativos, un gen se considera esencial cuando su pérdida de función impide la supervivencia o el crecimiento del organismo en un contexto experimental determinado. Esta definición requiere una interpretación cuidadosa: la esencialidad no es una propiedad fija del gen, sino que puede variar con el fondo genético, el medio de cultivo, el estado fisiológico y las presiones ambientales \cite{EsencialidadRancati2018,EsencialidadHogan2022}.

Los estudios de mutagénesis saturada con transposones permiten estimar, a escala genómica, qué regiones toleran inserciones y cuáles parecen indispensables en una condición dada. Complementariamente, la interferencia CRISPR puede reducir la expresión de genes esenciales y revelar vulnerabilidades que no son accesibles mediante deleción completa. En \textit{A. baumannii}, estas aproximaciones han identificado conjuntos de genes necesarios para la persistencia durante la infección, el crecimiento \textit{in vitro} y la respuesta a antibióticos \cite{EsencialidadWang2014,EsencialidadBai2021,EsencialidadWard2023}. No obstante, las diferencias entre cepas, bibliotecas, umbrales analíticos y condiciones ensayadas impiden interpretar una clasificación individual como evidencia universal. Por esa razón, en esta tesis la esencialidad se utiliza como una capa de evidencia que debe evaluarse junto con la conservación y los restantes atributos de cada candidato.

\section{Integración de evidencia y validación estructural}

La disponibilidad de genomas, ensayos funcionales y recursos de anotación posibilita establecer una estrategia reproducible de priorización. En primer lugar, la conservación en el pangenoma permite descartar candidatos de presencia restringida. Luego, la evidencia de esencialidad contribuye a estimar la vulnerabilidad funcional; la anotación de funciones y rutas metabólicas facilita interpretar el papel celular de cada proteína; y la búsqueda de homólogos humanos aporta un criterio inicial de selectividad. Finalmente, el análisis de estructuras proteicas y de cavidades de unión permite estimar si una proteína contiene sitios plausibles para el reconocimiento de ligandos y, por ende, si justifica una evaluación posterior de su ligabilidad \cite{Bosch2024}.

Este enfoque no pretende demostrar por sí solo que una proteína sea una diana clínicamente validada ni reemplazar la validación experimental. Su contribución es reducir de manera explícita y trazable el espacio de candidatos, identificar las incertidumbres asociadas a cada evidencia y seleccionar proteínas para estudios posteriores. En particular, combinar conservación, esencialidad y propiedades estructurales evita priorizar exclusivamente genes indispensables pero poco representativos de la diversidad clínica, o proteínas conservadas cuya intervención farmacológica resulte poco factible.

\section{Objetivos}

El objetivo general de esta tesis es desarrollar e implementar un flujo computacional integrador para priorizar potenciales dianas terapéuticas conservadas en cepas clínicas de \textit{Acinetobacter baumannii}, combinando evidencia genómica, funcional y estructural.

De manera específica, se propone analizar el pangenoma de una colección de cepas clínicas para identificar genes conservados; integrar evidencia publicada de esencialidad genética y contextualizar sus limitaciones; evaluar la anotación funcional, la participación en rutas metabólicas y la presencia de homólogos humanos de los candidatos; caracterizar propiedades estructurales y de ligabilidad de las proteínas mejor posicionadas; y generar un ranking interpretable y documentado que reúna la evidencia que sostiene la priorización de cada diana.
